\section{Methodology}
\label{sec:method} 

Figure~\ref{fig:LoopMTP_arch} contrasts our proposed method, \textsc{LoopMTP}, with standard looped transformers and multi-token prediction. \textsc{LoopMTP} addresses representation overwriting and latent overthinking, and undifferentiated computation through an MTP-guided looped block (Section~\ref{sec:method:cond}), learned state aggregation (Section~\ref{sec:method:agg}), and a soft-alignment auxiliary objective (Section~\ref{sec:method:loss}).





\subsection{Notation}
In the following, we build upon the notation of \citet{saunshi2025reasoning}. Let $f_\theta$ denote a stack of $L$ transformer blocks with shared
parameters $\theta$, and $T$ denote the number of times this stack is applied.
Given an input token sequence $\mathbf{u} {=} (u_1, \dots, u_S) \in \mathcal{V}^S$, where $\mathcal{V}$ denotes the vocabulary and $S$ the sequence length, we first embed the tokens:
\setlength\abovedisplayskip{5pt}
\setlength\belowdisplayskip{6pt}
\begin{equation}
  \mathbf{x}^{(0)} \;=\; \mathrm{Embed}(\mathbf{u}) \in \mathbb{R}^{S \times d},
\end{equation}
and then apply $f_\theta$ recursively for $t = 1, \dots, T$:
\begin{equation}
  \mathbf{x}^{(t)} \;=\; f_\theta\!\left(\mathbf{x}^{(t-1)}, \mathbf{x}^{(0)}\,,\, t\right).
  \label{eq:loop}
\end{equation}
The explicit dependence on $\mathbf{x}^{(0)}$ and the iteration index $t$
is discussed in Section~\ref{sec:method:cond}. 




\begin{figure*}[t]
    \centering
    \includegraphics[
        width=0.92\linewidth,
        trim=0cm 0pt 0cm 8pt, %
        clip
    ]{figures/rerendered/fig1.pdf}
    \vspace{-8pt}
    \caption{\textbf{Architecture and mechanism overview.} 
    (a)~A standard looped transformer block applied $T$ times where next-token supervision is applied only at the very end.
    (b)~A standard multi-token prediction (MTP) model requiring full-vocabulary projections.
    (c)~\textbf{\textsc{LoopMTP}} (ours): per-iteration outputs are aligned in the latent space with future token embeddings (soft MTP) and combined via a gating mechanism. Illustrated for a model with $T{=}3$ loops.}
    \vspace{-7pt}
    \label{fig:LoopMTP_arch}
\end{figure*}





\subsection{MTP-Guided Looped Block}
\label{sec:method:cond}
Our MTP-guided looped block $f_{\theta}$ comprises all $L$ transformer blocks, i.e., we loop over the entire model.
Each application of $f_\theta$ at iteration $t$ takes as input the previous
iteration's output $\mathbf{x}^{(t-1)}$ and the fixed token embeddings
$\mathbf{x}^{(0)}$, and outputs $\mathbf{x}^{(t)}$. We construct the
input to $f_\theta$ via three operations: (i) a normalized iteration-index
embedding, (ii) per-iteration normalization, and (iii) a concatenation-and-projection step.

\paragraph{Iteration-index embedding and per-iteration normalization.}
Following \citet{frey2026adaptive}, we explicitly signal the current iteration by appending the normalized iteration index to token vectors.
We then normalize hidden states with iteration-specific LayerNorms $\mathrm{LN}^{\text{prev}}_{t}$ and obtain token embeddings using
$\mathrm{LN}^{\text{tok}}$:
\begin{equation}
  \mathbf{h}^{(t-1)} \;{=}\; \mathrm{LN}^{\text{prev}}_{t}\!\left(\big[\, \frac{t{-}1}{T} \;\Vert\; \mathbf{x}^{(t-1)} \,\big]\right) {\in } \mathbb{R}^{S \times (d+1)}, 
\end{equation}
\begin{equation}
  \mathbf{e}   \;=\; \mathrm{LN}^{\text{tok}}\!\left(\mathbf{x}^{(0)}\right) \in \mathbb{R}^{S \times d},
\end{equation}
where $\Vert$ denotes feature-axis concatenation.
Per-iteration norms are crucial, letting each loop operate on inputs with iteration-specific statistics.

\paragraph{Input fusion.}
The previous-iteration state and the re-normalized token embedding are merged
by concatenation along the feature axis followed by a linear projection $\mathrm{P}$:

\begin{equation}
  \mathbf{y}^{(t)} \;=\; \mathrm{P}\!\left(\big[\, \mathbf{e} \;\Vert\; \mathbf{h}^{(t-1)} \,\big]\right) \in \mathbb{R}^{S \times d}.
\end{equation}

The shared backbone is then applied:
\begin{equation}
  \mathbf{x}^{(t)} \;=\; f_\theta\!\left(\mathbf{y}^{(t)}\right)
                    \;=\; \mathrm{Layer}_L \circ \cdots \circ \mathrm{Layer}_1\!\left(\mathbf{y}^{(t)}\right).
\end{equation}


\paragraph{Looped layer-norm scaling (Loop-LNS).}
To preserve stability under a variable number of unrollings, we adapt layer-norm scaling (LNS) \citep{sun2026curse} to the looped setting. Inside each layer $\ell$, the input $\mathbf{x}$ is first normalized (i.e.~pre-norm transformer) and rescaled by a fixed factor $\frac{1}{T}$:
\begin{align}
  \mathbf{x} \;=\; \mathbf{x} + \mathrm{Attn}\!\left(\tfrac{1}{T}\,\mathrm{RMSNorm}(\mathbf{x})\right),
  \\
  \mathbf{x} \;=\; \mathbf{x} + \mathrm{FFN}\!\left(\tfrac{1}{T}\,\mathrm{RMSNorm}(\mathbf{x})\right).
\end{align}
Unlike standard LNS \citep{sun2026curse}, which scales by $\frac{1}{\sqrt{n}}$, where $n$ is the running depth, our \emph{Loop-LNS} variant uses a fixed factor of $\frac{1}{T}$ for all 
effective sub-blocks, reflecting the fact that the backbone is reused $T$ times rather than deepened.

We chose the fixed per-iteration factor $\frac{1}{T}$ over running-depth alternatives for a specific reason: a progressive counter that scales layer~$\ell$ in iteration~$t$ by $\frac{1}{(t{-}1) \cdot L + \ell}$  yields tiny factors for later iterations of deep unrollings, which we found empirically neutralizes their gradient contribution. 
The constant $\frac{1}{T}$ balances residual-stream stability with sufficient signal propagation to all iterations.


\paragraph{Per-Iteration Representation Alignment with MTP. }
To exploit the per-iteration hidden states $\{\mathbf{x}^{(t)}\}_{t=2}^{T}$, we encourage each iteration's hidden state to point toward the output embedding of the token it should anticipate. 
For iteration $t {\geq} 2$, position $i$ is aligned to the output embedding of token $u_{i+t}$, so that the $t$-th iteration is responsible for anticipating the token $t$ steps ahead.
The first iteration $t{=}1$ is reserved as an unconstrained representation (i.e.~no supervision): the model is free to learn what kind of information to encode here, so that this representation is rich enough for subsequent iterations to read out multiple future tokens from it.


\subsection{Aggregation}
\label{sec:method:agg}

Rather than discarding the intermediate iterations, we aggregate
$\{\mathbf{x}^{(t)}\}_{t=1}^{T}$ via a token-wise weighted sum with a
\emph{shared, content-conditional} gate. %
Let $W_g \in \mathbb{R}^{d \times d}$
be a single linear gate shared across all iterations and let
$\boldsymbol{\beta} = (\beta_1, \dots, \beta_T) \in \mathbb{R}^{T}$
be per-iteration scalar bias parameters. For each iteration $t$ and position
$i \in \{1,\dots,S\}$, we compute an unnormalized gate:
\begin{equation}
  g_i^{(t)} = \text{softplus}\!\Big(W_g\, \mathbf{x}_i^{(t)} + \beta_t \cdot \mathbf{1}_d\Big)
  \;\in\; \mathbb{R}^d,
\end{equation}
where $\mathbf{1}_d \in \mathbb{R}^d$ is the all-ones vector. We then normalize gate values across iterations:
\begin{equation}
  \tilde{\mathbf{g}}^{(t)}_{i} {=} \frac{\mathbf{g}^{(t)}_{i}}{\sum_{s=1}^{T} \mathbf{g}^{(s)}_{i} + \varepsilon},
\end{equation}
where $\varepsilon = 10^{-8}$ to avoid division by zero. 
Finally, the aggregation happens as follows:
\begin{equation}
    \mathbf{z}_{i} \;=\; \sum_{t=1}^{T} \tilde{\mathbf{g}}^{(t)}_{i} \odot \mathbf{x}^{(t)}_{i},
  \label{eq:ws}
\end{equation}
where $\odot$ is the elementwise product. 
The gating mechanism itself introduces a linear layer of size $d {\times} d$ and $T$ learnable scalar bias terms, amounting to ${<}0.5\%$ additional parameters for our models.

Similarly to \citet{frey2026adaptive}, we found the gate initialization important: 
the exact initial value matters less than how the model behaves at the beginning of training.
The first bias term is initialized to a moderately positive value, while all subsequent bias terms start at a substantially negative value. This provides a strong prior toward the first iteration, i.e., acting as a vanilla non-looped transformer at first, while leaving later iterations free to take over as training progresses.










\subsection{Training Objective}
\label{sec:method:loss}

The total loss combines three terms: the main next-token cross-entropy, a hidden-state alignment loss, and a
ponder regularizer.

\paragraph{Main NTP loss.} The next-token prediction (NTP) loss, commonly used as the main pretraining objective for LLMs, is applied to the output of the language-model head on the normalized aggregated representation $\mathbf{z}$ (denoted by $\mathbf{p}_i^{\text{agg}}$):
\begin{equation}
  \mathcal{L}_{\text{NTP}} \;=\; -\frac{1}{S-1}\sum_{i=1}^{S-1} \log \mathbf{p}_i^{\text{agg}}\!\left[u_{i+1}\right].
\end{equation}


\paragraph{Hidden-state soft MTP alignment.}
We regularize the per-iteration hidden states to align with the embeddings of the tokens they are tasked with predicting. Let $E \in \mathbb{R}^{|\mathcal{V}| \times d}$ denote the output (unembedding) matrix, treated as a fixed target via stop-gradient $\mathrm{sg}[\cdot]$. For $t = 2, \dots, T$, the per-step hidden-state alignment loss is:
\begin{equation}
  \mathcal{L}^{(t)}_{\text{align}} \;=\; \frac{1}{S - t}\sum_{i=1}^{S - t}\!\left(1 - \cos\!\big(\mathbf{x}^{(t)}_{i},\, \mathrm{sg}[E_{u_{i+t}}]\big)\right).
\end{equation}
The final alignment loss averages these:
\begin{equation}
  \mathcal{L}_{\text{align}} \;=\; \frac{1}{T-1}\sum_{t=2}^{T}\mathcal{L}^{(t)}_{\text{align}}.
\end{equation}


Notably, in this case we align the per-iteration \emph{hidden state} with the future token \emph{embedding}, in contrast to $\mathcal{L}_\text{NTP}$, where the alignment is over the \emph{distribution}. By matching embeddings rather than full output distributions, our objective imposes a softer constraint, which is why we refer to it as \emph{soft} MTP alignment. 
We emphasize that our alignment approach does not predict multiple future tokens
in the traditional sense: each iteration's hidden state is steered toward a future
token's embedding via cosine similarity, acting as a \emph{representational regularizer}
rather than a distributional constraint.  We retain the term ``multi-token
prediction'' to reflect the lookahead structure of the supervisory signal, while
noting that this cosine-based formulation avoids the vocabulary-sized projection that makes standard MTP expensive---a property that we believe contributes to its
effectiveness at small scale.
Importantly, computing $\mathcal{L}_\text{align}$ requires only a cosine similarity between each iteration's output and its corresponding target, adding negligible overhead.



\paragraph{Ponder regularizer.}
The aggregator's gates induce a per-token distribution
$G_i(t) {=} \tfrac{1}{d}\sum_{k} \tilde{\mathbf{g}}^{(t)}_{i,k}$ over the
$T$ iterations. 
Following \citet{zhu2025scaling}, we pull this distribution toward the uniform
prior $Q {=} (1/T, \dots, 1/T)$ over the $T$
iterations via the Kullback--Leibler (KL) divergence:
\begin{equation}
  \mathcal{L}_{\text{ponder}} \;=\; \frac{1}{S}\sum_{i=1}^{S} \mathrm{KL}\!\left(G_i \,\Vert\, Q\right).
\end{equation}


\paragraph{Final loss.}
The final loss is calculated as the weighted sum of next token prediction loss, alignment loss, and the ponder loss as follows:
\begin{equation}    
    \mathcal{L} = \mathcal{L}_{\text{NTP}}
      + \lambda_{\text{align}}\, \mathcal{L}_{\text{align}}
      + \lambda_{\text{ponder}}\, \mathcal{L}_{\text{ponder}}
  \label{eq:total}
\end{equation}
